%%%%%%%%%%%%%%%%%%%%%%%%%%%%%%%%%%%%%%%%%%%%%%%%%%%%%%%%%%%%%%%%%%%%%%%%%%%%%%%%
%%
\cleardoubleoddpage%  Make sure to start each chapter on a new odd page
\chapter{Background}
\label{ch:background}


\section{Geographic Atrophy and OCT Imaging}
\label{sec:background:ga-oct}

This section provides the imaging-side background needed to motivate
the state representation used throughout the remainder of the thesis.
The clinical framing established in Chapter~\ref{ch:introduction} is
taken as read; the focus here is on \emph{which part of the eye} GA
affects, on \emph{how} that part is imaged, on \emph{what is visible}
in an OCT volume of a GA eye, and on \emph{which derived quantities}
feed into the eleven-channel state tensor introduced
in~\S\ref{sec:data:state}. Because the rest of this thesis is written
for a primarily machine-learning audience, the present section is
deliberately self-contained and aims to remain accessible to readers
without prior clinical training:~\S\ref{sec:background:ga-oct:anatomy}
introduces the small set of anatomical terms that recur throughout the
work, and the subsequent subsections build the imaging substrate on
top of that vocabulary.

\subsection{A short anatomy primer}
\label{sec:background:ga-oct:anatomy}

% TODO: produce Figure fig:bg:eye-anatomy --- a two-panel anatomy
% reference for ML readers without clinical background:
% (left) sagittal cross-section of the human eye labelling the
% principal structures from front to back: cornea, lens, vitreous
% body, retina, fovea, optic nerve head (ONH) / optic disc, optic
% nerve, choroid, sclera; with an explicit annotation of the macular
% region as a small disc-shaped area on the retina centred near the
% posterior pole. (right) en-face fundus view of the posterior pole
% looking at the inside of the back of the eye, labelling: macula
% (~6 mm-wide central region), fovea (central pit), foveola (the
% innermost ~0.35 mm of the fovea), parafovea (~0.5--1.5 mm
% eccentric ring), perifovea (~1.5--3 mm eccentric ring), and ONH /
% optic disc (~4 mm nasal to the fovea). The 6 x 6 mm OCT scanning
% window used by the thesis cohort can be overlaid on the fundus
% panel as a dashed square centred on the fovea, to make the link
% between the anatomical reference frame and the modelling grid
% explicit.
\begin{figure}[t]
  \centering
  \fbox{\parbox[c][6cm][c]{0.9\textwidth}{\centering\itshape Placeholder
    for Figure~\ref{fig:bg:eye-anatomy}: sagittal cross-section of the
    eye (left) and en-face view of the fundus (right) labelling the
    anatomical terms used throughout this thesis. See accompanying
    \texttt{\% TODO} for the intended content.}}
  \caption[Anatomy of the eye and the macular region.]{%
    Anatomy reference for the terms used throughout this thesis.
    \emph{Left:} sagittal (side-view) cross-section of the eye. Light
    enters through the cornea and lens at the front, traverses the
    vitreous body, and falls on the retina at the back; the
    \emph{macula} is a small disc-shaped specialised region of the
    retina near the posterior pole, and the \emph{fovea} is its
    centre, where photoreceptor density and visual acuity peak. The
    optic-nerve head (ONH, also called the optic disc) is the point
    at which the optic nerve exits the eye.
    \emph{Right:} en-face view of the fundus, i.e.\ the inside of the
    back of the eye seen along the optical axis. The fovea sits
    roughly at the centre of the macula; the macula itself is
    sub-divided into the foveola (innermost), parafovea, and
    perifovea, in concentric rings of increasing radius. The dashed
    square indicates the
    $6 \times 6$\,mm fovea-centred OCT scanning window used in this
    thesis (\S\ref{sec:background:ga-oct:muw}).}
  \label{fig:bg:eye-anatomy}
\end{figure}

The eye, sketched in Figure~\ref{fig:bg:eye-anatomy} (left), can be
thought of as a biological camera. Light enters at the front through
the \emph{cornea} and \emph{lens}, traverses the gel-filled vitreous
body, and is focused onto the \emph{retina} --- a thin
neural-tissue lining at the back of the eye that converts the
incoming light pattern into electrical signals. Those signals leave
the eye via the \emph{optic nerve}, which exits the retina at a
specific anatomical landmark called the \emph{optic-nerve head} (ONH,
also called the \emph{optic disc}). Behind the retina, two further
layers --- the \emph{choroid} (a richly vascularised tissue, of which
the \emph{choriocapillaris} is the innermost capillary bed) and the
\emph{sclera} (the white, mechanically rigid outer shell) --- support
and supply the retina from beneath.

The retina is not uniform across its surface. A small disc-shaped
region near the posterior pole, the \emph{macula}, is specialised for
high-acuity central vision, while the rest of the retina handles
peripheral vision at lower resolution. The macula has a diameter of
roughly five to six millimetres and is the only part of the retina
that GA ever affects clinically; for that reason, the OCT scanning
window used throughout the thesis is centred on the macula and
covers a $6 \times 6$\,mm field
(\S\ref{sec:background:ga-oct:muw}). Within the macula, four
concentric anatomical zones are conventionally distinguished, in
order of increasing radius from the centre: the \emph{foveola}, the
\emph{fovea}, the \emph{parafovea}, and the \emph{perifovea}
(Figure~\ref{fig:bg:eye-anatomy}, right). The fovea is the small
central pit at which photoreceptor density and visual acuity reach
their peak, and the term ``foveal'' refers to anything in or close to
that pit; ``parafoveal'' and ``perifoveal'' describe progressively
more eccentric annular regions. Two qualifiers used in the GA
literature follow directly from this geography: a lesion is called
\emph{subfoveal} if it overlaps the fovea (i.e.\ it directly threatens
high-acuity vision) and \emph{extrafoveal} otherwise. As discussed
in~\S\ref{sec:introduction:problem}, whether a lesion is subfoveal or
extrafoveal materially changes its clinical significance even if its
total area is the same.

A second piece of anatomy is internal rather than positional: the
retina is itself a \emph{layered} tissue, with several distinct
strata stacked vertically from the front (closest to the vitreous
body) to the back (closest to the choroid). The principal layers, in
order from inner to outer, are the inner limiting membrane, the
retinal nerve fibre layer, the ganglion-cell and inner-plexiform
layers, the inner nuclear and outer plexiform layers, the outer
nuclear layer, the photoreceptor inner- and outer-segment bands, the
retinal pigment epithelium (RPE), and Bruch's membrane;
Figure~\ref{fig:bg:retinal-anatomy} (left) gives a labelled schematic.
The terminology in this list takes some getting used to, but for the
purposes of this thesis it is sufficient to remember the structural
fact that the retina is layered, that GA destroys a specific
\emph{outer}-retinal stack of three of those layers
(\S\ref{sec:background:ga-oct:disease}), and that the surrounding
layer geometry --- ten boundary depth maps in total --- is what feeds
the model alongside the lesion mask
(\S\ref{sec:background:ga-oct:state-rep}).

The remaining subsections build on this vocabulary. \S\ref{sec:background:ga-oct:disease}
describes how GA appears on OCT; \S\ref{sec:background:ga-oct:oct}
explains what an OCT volume actually is and how it is acquired;
\S\ref{sec:background:ga-oct:enface} describes the projection from
3D OCT volume to the 2D en-face grid that the model operates on;
\S\ref{sec:background:ga-oct:state-rep} defines the eleven channels of
the state tensor; and \S\ref{sec:background:ga-oct:muw} describes the
MUW longitudinal cohort.

\subsection{Geographic Atrophy on OCT}
\label{sec:background:ga-oct:disease}

% TODO: produce Figure fig:bg:retinal-anatomy --- a two-panel
% schematic for readers without a clinical background:
% (left) labelled cross-section through a healthy macula showing the
% principal retinal layers from ILM at the top through RNFL, GCL+IPL,
% INL+OPL, ONL, photoreceptor IS/OS bands (including the ellipsoid
% zone), RPE, and Bruch's membrane, with the choriocapillaris
% immediately beneath; (right) a representative OCT B-scan through a
% GA-affected macula highlighting the dropout of the outer-retinal
% layers within the lesion and the corresponding choroidal
% hypertransmission signature beneath. Mirror the channel ordering of
% the eleven-channel state tensor used in the thesis.
\begin{figure}[t]
  \centering
  \fbox{\parbox[c][6cm][c]{0.9\textwidth}{\centering\itshape Placeholder for
    Figure~\ref{fig:bg:retinal-anatomy}: schematic cross-section of a
    healthy macula (left) versus a GA-affected macula on OCT (right).
    See accompanying \texttt{\% TODO} for the intended content.}}
  \caption[Retinal anatomy and the OCT signature of geographic atrophy.]{%
    Anatomy of the macular retina and the OCT signature of geographic
    atrophy. \emph{Left}: schematic cross-section through a healthy
    macula, labelling the principal retinal layers from the inner
    limiting membrane (ILM) at the top --- through the retinal nerve
    fibre layer, the ganglion-cell and inner-plexiform layers, the inner
    nuclear and outer plexiform layers, the outer nuclear layer, and the
    photoreceptor inner- and outer-segment bands including the ellipsoid
    zone --- down to the retinal pigment epithelium (RPE) and Bruch's
    membrane (BM), with the choriocapillaris immediately beneath. The
    same set of anatomical strata corresponds to the ten layer-boundary
    channels of the eleven-channel state tensor introduced in
    \S\ref{sec:background:ga-oct:state-rep}. \emph{Right}: a
    representative OCT B-scan through a GA-affected macula. Inside the
    lesion the outer-retinal complex is missing and a brighter signal
    appears in the choroid beneath the lesion (the
    \emph{choroidal hypertransmission} signature, item~(ii) in the
    text below).}
  \label{fig:bg:retinal-anatomy}
\end{figure}

The retina is a thin, layered tissue at the back of the eye that
converts incoming light into neural signals;
Figure~\ref{fig:bg:retinal-anatomy} (left) sketches its principal
anatomical strata in the macula. Geographic atrophy is the disease in
which a stack of three tightly coupled \emph{outer}-retinal layers ---
the retinal pigment epithelium (RPE), the photoreceptor outer segments
that sit just above the RPE, and the choriocapillaris (a dense capillary
bed) immediately below the RPE --- die together in a localised patch of
the macula \citep{Boopathiraj2024,Boyer2017}. These three layers are
anatomically interdependent: the RPE supports the photoreceptors, and
the choriocapillaris supplies both. In other words, atrophy of one
layer entrains the degeneration of the others, so the resulting lesion
appears on imaging as a sharply delineated patch in which the entire
outer-retinal complex is missing rather than as the loss of a single
isolated layer.

What this end-stage lesion looks like on OCT has been formalised by an
international consensus group, the Classification of Atrophy Meetings
(CAM), under the term \emph{complete RPE and outer retinal atrophy}
(cRORA) \citep{Vallino2024}. cRORA is identified by three concurrent
OCT features: \emph{(a)} a region of choroidal hypertransmission of at
least 250\,\textmu{}m in diameter; \emph{(b)} co-located attenuation or
disruption of the RPE band of at least 250\,\textmu{}m; and \emph{(c)}
overlying photoreceptor degeneration --- with the additional condition
that no abnormal blood-vessel growth (\emph{macular neovascularisation},
which is the wet form of AMD) is present in the same region.
% TODO: cite Sadda2018 (CAM Report 3, "Consensus Definition for Atrophy
% Associated with Age-Related Macular Degeneration on OCT", Ophthalmology
% 125:537-548) once added to references.bib -- this is the primary
% source for the cRORA criteria attributed to Vallino2024 here.
Put differently, cRORA is the OCT-readable definition of an established
dry-AMD lesion, and \emph{GA is precisely the subset of cRORA that
arises in eyes without macular neovascularisation}. The CAM framework
also recognises a precursor stage, \emph{iRORA} (incomplete RPE and
outer retinal atrophy), which is essentially the same lesion at a
smaller spatial scale, i.e.\ the same OCT pattern below the
250\,\textmu{}m size threshold; iRORA progresses to cRORA in
approximately 93\,\% of cases within twenty-four months, with a median
conversion time of fourteen months \citep{Vallino2024}. The thesis
state tensor encodes only the cRORA endpoint, as channel~0; iRORA is
implicit in the surrounding layer channels but is not represented as a
separate disease label, a scoping decision discussed
in~\S\ref{sec:discussion:limitations}.

The principal direct OCT signatures of an established GA lesion ---
the visual fingerprints by which the disease is identified on a B-scan
--- are the following \citep{Boopathiraj2024,Yehoshua2011,Vallino2024}.
\emph{(i)} The outer-retinal layers are missing: the photoreceptor
inner-segment / outer-segment junction (the \emph{ellipsoid zone}) and
the RPE band drop out of the B-scan, leaving a structurally collapsed
column where intact retina used to be.
\emph{(ii)} \emph{Choroidal hypertransmission} appears beneath the
lesion: because pigmented RPE no longer absorbs the probe beam, light
penetrates deeper into the choroid, and the column-summed signal is
higher than in the surrounding intact tissue. This is the same optical
effect that makes GA appear bright on summed-voxel en-face projections,
discussed in~\S\ref{sec:background:ga-oct:enface}.
\emph{(iii)} \emph{RPE--Bruch's-membrane (BM) thickening} is
concentrated immediately around the lesion margin --- essentially, the
narrow ring of tissue just outside the lesion is measurably thicker
than healthy tissue elsewhere on the retina --- and is attributed to
basal-laminar deposits accumulated in the diseased tissue; the
thickness of this rim correlates strongly with the eye's subsequent
annual square-root lesion enlargement rate \citep{Chu2022}.
The first two signatures define the lesion mask used as channel~0 of
the state tensor; the third is a boundary-localised predictive
biomarker that motivates encoding the surrounding layer geometry
rather than the mask alone.

\subsection{Optical Coherence Tomography of the Macula}
\label{sec:background:ga-oct:oct}

% TODO: cite Huang1991 (Huang et al., Science 254:1178-1181, "Optical
% Coherence Tomography") once added to references.bib -- the seminal
% paper introducing OCT, conventionally cited as the methodological
% origin of the modality.
% TODO: cite Wojtkowski2002 / Drexler2008 (or an equivalent SD-OCT
% review) once added to references.bib -- the principle behind
% spectral-domain OCT (the modality used to acquire the MUW cohort)
% conventionally rests on these works.

% TODO: produce Figure fig:bg:oct-acquisition --- a four-panel
% schematic showing how an OCT volume is built up from one-dimensional
% depth profiles, intended to anchor the geometric vocabulary
% (A-scan, B-scan, volume, en-face) for ML readers without imaging
% background:
% (a) a sketch of the eye with a single near-infrared probe beam
%     entering the cornea and reaching the retina; an inset graph
%     shows the resulting depth-versus-reflectivity profile (= one
%     A-scan), with characteristic peaks at the principal retinal
%     layer interfaces.
% (b) the same sketch with the beam scanned laterally along one
%     axis, producing a stack of adjacent A-scans which when
%     displayed side-by-side form a 2D cross-sectional image
%     (= one B-scan).
% (c) the beam now scanned in both lateral directions, producing
%     a 3D rectangular cuboid of voxels (= one OCT volume) with the
%     physical footprint explicitly annotated.
% (d) the volume projected (e.g. by summing along the axial
%     direction) onto its top face, yielding the 49 x 1024 en-face
%     image used by the thesis (= the en-face projection of
%     Section~2.1.4). A representative en-face image with a GA
%     lesion visible as a brighter region can be shown next to the
%     volume to make the projection step concrete.
\begin{figure}[t]
  \centering
  \fbox{\parbox[c][6cm][c]{0.9\textwidth}{\centering\itshape Placeholder
    for Figure~\ref{fig:bg:oct-acquisition}: four-panel schematic of
    OCT acquisition geometry, going A-scan $\to$ B-scan $\to$ volume
    $\to$ en-face projection. See accompanying \texttt{\% TODO} for
    the intended content.}}
  \caption[OCT acquisition geometry: A-scan, B-scan, volume,
  en-face projection.]{%
    OCT acquisition geometry, showing how a three-dimensional OCT
    volume is built up from one-dimensional depth profiles and how
    the two-dimensional en-face image used by the model relates to
    that volume.
    \emph{(a)} A single near-infrared probe beam enters the eye and
    is reflected back from the various retinal and choroidal layers
    along its path; the resulting depth-versus-reflectivity profile
    is one \emph{A-scan}, and the bright bands in that profile
    correspond to the principal layer interfaces.
    \emph{(b)} Scanning the beam laterally along one direction
    produces a sequence of A-scans whose side-by-side display forms
    a two-dimensional cross-sectional image of the retina, called a
    \emph{B-scan}.
    \emph{(c)} Scanning in both lateral directions yields a
    three-dimensional \emph{OCT volume} (the physical scan window of
    the MUW cohort spans approximately $6 \times 6$\,mm laterally).
    \emph{(d)} Projecting the volume along its axial direction onto
    the fundus plane yields the two-dimensional en-face image on
    which the modelling pipeline operates
    (\S\ref{sec:background:ga-oct:enface}); a GA lesion appears
    inside the en-face image as a region with a different signal
    intensity than the surrounding healthy retina.}
  \label{fig:bg:oct-acquisition}
\end{figure}

Optical Coherence Tomography (OCT) produces micrometre-resolution
cross-sectional images of the retina, in much the same way that
ultrasound produces them for soft tissue --- except that low-coherence
near-infrared light replaces sound. Concretely, the system shines a
weak probe beam into the eye and measures how strongly each
\emph{depth} along the path of the beam reflects light back. That
depth-versus-reflectivity profile is then displayed as a single column
of pixels in the output image, with intensity encoding reflectivity
and the vertical position of each pixel encoding depth. Resolving the
depth dimension requires a small piece of optical engineering, because
near-infrared light from a single short pulse cannot be timed
precisely enough on a per-layer basis with conventional electronics.
The standard solution is interferometric: the source light is split
into two beams (Figure~\ref{fig:bg:oct-acquisition}(a)), one of which
--- the \emph{sample arm} --- is sent into the eye and reflected back
from the various retinal layers, while the other --- the
\emph{reference arm} --- is sent to a known mirror and reflected back
unchanged. The two returning beams are then recombined and produce an
interference pattern, and the spectrum of that pattern (in
spectral-domain OCT) encodes the depth at which each reflection in
the eye occurred. Standard signal processing then recovers the
A-scan as a one-dimensional function of depth.

A single A-scan is one column of one image. Scanning the probe beam
laterally along one direction produces many adjacent A-scans, and
displaying them side by side yields a two-dimensional cross-section
through the retina, called a \emph{B-scan}
(Figure~\ref{fig:bg:oct-acquisition}(b)). Scanning in both lateral
directions yields a three-dimensional \emph{OCT volume}
(Figure~\ref{fig:bg:oct-acquisition}(c)), in which each voxel records
the reflectivity of one small block of tissue. The MUW cohort uses an
acquisition protocol that produces volumes with a physical footprint of
approximately $6 \times 6$\,mm laterally; how the varying native voxel
grids are reduced to the standard two-dimensional modelling grid used in
the rest of the thesis is the subject of~\S\ref{sec:background:ga-oct:enface}.

OCT is preferred over fundus autofluorescence (FAF), the other common
imaging modality for GA, for three reasons given by \citet{Vallino2024}: it
shows each retinal layer in cross-section and so permits layer-wise
quantitative analysis, it detects atrophy earlier in its natural history,
and it remains reliable near the fovea, where blue-light FAF struggles.
OCT-based area measurements are also reproducible and consistent with FAF:
\citet{Pilotto2015} report inter-reader intraclass correlation coefficients
of 0.97--0.99 for en-face OCT and find the resulting GA areas statistically
equivalent to blue-light and near-infrared FAF measurements in a small
cohort. The data used in this thesis are therefore OCT-derived.

The single OCT acquisition simultaneously yields all the disease
information that the state tensor needs to encode. The lesion itself is
visible on the B-scan as a region of outer-retinal layer dropout with
choroidal hypertransmission; the surrounding retinal layer geometry is
recoverable from the same volume by automatic segmentation; and the
choroidal slab beneath Bruch's membrane provides an additional view of
the GA extent through the hypertransmission signature
\citep{Vallino2024,Yehoshua2011}. No cross-modality fusion is required.

\subsection{From OCT volumes to en-face projections}
\label{sec:background:ga-oct:enface}

The full OCT volume of a GA eye is too large and too redundant to be
modelled at native voxel resolution. The clinical literature has
converged on a two-dimensional \emph{en-face projection} of the volume
onto the fundus plane as the canonical analysis substrate
\citep{Pilotto2015,Chu2022,Yehoshua2011,Vallino2024}, and the present
work adopts the same convention.

Two construction strategies appear in the literature. The
\emph{summed-voxel projection} of \citet{Yehoshua2011} integrates the
A-scan intensity along the full axial extent of the volume; in this
projection GA appears as a \emph{bright} region because the absence of
the pigmented RPE column lets choroidal signal contribute to the column
sum that would otherwise be absorbed. The \emph{slab segmentation}
strategy of \citet{Pilotto2015} instead extracts a thin axial slab of
the volume and projects only that slab onto the fundus plane; an outer
retinal slab placed just above Bruch's membrane visualises RPE and
photoreceptor-complex atrophy directly, whereas a sub-RPE slab placed
between 64 and 400\,\textmu{}m below Bruch's membrane visualises the
same lesion through its choroidal hypertransmission signature
\citep{Chu2022,Vallino2024}. The two strategies converge on the same
underlying lesion geometry.

The choice of en-face rather than B-scan as the modelling substrate is
motivated by a methodologically important caveat documented by
\citet{Vallino2024}: a lesion classified as iRORA on a single horizontal
B-scan can in fact be part of a larger cRORA lesion that is only fully
visible in the en-face slab. The en-face representation is therefore
the more complete view of the GA extent and is the substrate on which
quantitative growth-rate metrics, drug-response readouts, and
deep-learning progression predictors are reported in the recent
literature \citep{Chu2022,Vogl2021,Vallino2024}. By construction, the
2D en-face grid also reduces a per-step three-dimensional forward
problem to a per-step two-dimensional one, which matters for the
tractability of the graph-neural-network architectures introduced in
the subsequent chapters.

\subsection{The eleven-channel state representation}
\label{sec:background:ga-oct:state-rep}

This subsection describes \emph{what} the eleven channels of the state
tensor are, in structural terms, and \emph{why} retaining all of them
is clinically motivated. The operational details of how the channels
are produced from the raw OCT volumes in the MUW cohort ---
segmentation pipeline, grading protocol, software stack --- are
deferred to~\S\ref{sec:data:dataset} and are not relied upon
in the rest of this chapter.
% TODO: confirm with supervisor the exact provenance of channel 0
% (manual clinical grading vs trained segmentation network vs
% combination) and of channels 1-10 (which layer-segmentation
% algorithm, which boundary list, which post-processing) before
% drafting Chapter 3.1; the present subsection is written to be
% agnostic of those details.

The first channel of the state tensor is a binary \emph{GA mask} on the
en-face grid: an indicator function that takes value one inside the
cRORA region defined by the CAM criteria
of~\S\ref{sec:background:ga-oct:disease} and zero outside it. The only
properties of channel~0 that this chapter relies upon are that it is
binary, that it is aligned to the same en-face grid as the layer
channels, and that the supported region of the mask coincides with the
GA lesion as it would be identified by a trained reader applying the
cRORA criteria.

The remaining ten channels are \emph{retinal layer boundary depth
maps}: at every pixel of the en-face grid, the value of channel
$c \in \{1, \dots, 10\}$ is the axial position of one anatomical layer
boundary in the underlying OCT B-scan. Concretely, fixing a pixel of
the en-face grid selects a single A-scan in the OCT volume, and channel
$c$ at that pixel records how deep along that A-scan the $c$-th
anatomical layer-boundary surface lies. Stacking all ten channels
therefore reconstructs the full vertical retinal architecture of the
eye on the same horizontal grid as the lesion mask. The specific list
of boundaries used for the MUW cohort is given
in~\S\ref{sec:data:state}; collectively the boundaries delimit the
principal anatomical strata of the retina, from the inner limiting
membrane at the top, through the retinal nerve fibre layer, the
ganglion-cell and inner-plexiform layers, the inner nuclear and outer
plexiform layers, the outer nuclear layer, the photoreceptor inner-
and outer-segment bands (including the ellipsoid zone), the retinal
pigment epithelium, and Bruch's membrane at the bottom (cf.\
Figure~\ref{fig:bg:retinal-anatomy}).

Two technical caveats are inherent to standard SD-OCT-derived layer
segmentation, regardless of the specific algorithm used to produce the
boundary maps, and are explicitly acknowledged here and again
in~\S\ref{sec:discussion:limitations}. First, standard SD-OCT cannot
separate the outer nuclear layer from the Henle fibre layer
\citep{Vogl2021}, so what is referred to as the ``ONL'' boundary in any
SD-OCT-derived state tensor is effectively a combined
ONL\,+\,HFL boundary. Second, the hyperreflective complex that
aggregates the RPE, the ellipsoid zone, the interdigitation zone, and
the outer photoreceptor segments sits at or below the axial resolution
limit of SD-OCT and is conventionally treated as a single composite
band rather than resolved into its sub-bands. Both caveats apply to the
state tensor used throughout the thesis.

The clinical motivation for retaining all ten layer-boundary channels
alongside the binary mask is that documented prognostic biomarkers for
GA progression live throughout the full vertical retinal column, not
only inside the lesion itself. RPE--Bruch's-membrane thickening at the
0--300\,\textmu{}m rim around the GA boundary correlates with annual
square-root enlargement rate at $r = 0.595$ \citep{Chu2022}; outer
retinal layer thinning, in particular of the outer nuclear layer and
of the RPE-photoreceptor complex, accelerates in the months preceding
macular-atrophy conversion \citep{Vogl2021}; and the inner nuclear
layer in the retina immediately adjacent to a GA lesion correlates
with subsequent yearly progression rate at Spearman $r = 0.48$
\citep{Ebneter2016}. \citet{Vallino2024} further report that
RPE-to-Bruch's-membrane distance combined with outer-retinal-layer
thickness jointly predicts approximately 62\,\% of the annual GA
progression variance, while either feature alone explains less than
ten per cent. The layer channels are therefore not auxiliary context;
they carry rate-of-change information that the mask alone cannot
provide. \S\ref{sec:method:multichannel} returns to how this
multi-channel state propagates through the network.

\subsection{The MUW longitudinal GA cohort}
\label{sec:background:ga-oct:muw}

The dataset used throughout the thesis is the GA cohort of the Medical
University of Vienna (MUW), a longitudinal collection of macula-centred
SD-OCT volumes acquired in routine clinical practice. The cohort comprises
553 Heidelberg SD-OCT volume scans over 75 eyes. Lateral pixel counts
vary per eye across the 553 visits, ranging from 38 to 66 B-scans by 961
to 1719 A-scans, giving 67 distinct native shapes in total. No visit has
a native dimension of $49 \times 1024$. Instead, $49 \times 1024$ is the
modelling grid, reached by a symmetric centre-crop or zero-pad operation,
and never by resampling (\S\ref{sec:data:spatial}). Pixel spacing is
treated as constant across the cohort at approximately
$[0.12118,\, 0.003867,\, 0.00568]$\,mm in the $(x, z, y)$ directions,
because the scans were resampled upstream to one gold-standard spacing.
The resulting fixed physical window is approximately
$5.94 \times 5.82$\,mm.
% TODO: check the DICOM spacing fields to confirm the constant-spacing assumption.
This constant-spacing assumption has only been checked against the
cohort's own recorded transform files, where it holds to a tolerance of
about two per cent. The ratio between the two in-plane spacings is
roughly 21:1. This extreme anisotropy is a direct consequence of the
clinical acquisition protocol and is reflected in several architectural
choices in Chapter~\ref{ch:method}.

The 6\,mm window is a standard in the literature. \citet{Pilotto2015}
explicitly exclude eyes whose GA extends beyond a $6 \times 6$\,mm
fovea-centred square, observing that approximately ninety-nine per cent
of GA pathology occurs within a central three-millimetre radius from
the macula. The same window is used by \citet{Yehoshua2011}
($200 \times 200$ raster at 30\,\textmu{}m / pixel lateral) and by
\citet{Chu2022} (Zeiss PLEX Elite 9000 swept-source OCT, $6 \times 6$\,mm scan
window with 1536 axial pixels and 500 A-lines per B-scan), and it
serves as the de-facto analysis area for GA OCT studies. However,
measured on the present cohort, this window is not lossless. It cuts
real lesion area in 27.3 per cent of visits. The cut removes more than
5 per cent of the lesion in 6.0 per cent of visits, reaching a
worst-case loss of 25.7 per cent. Furthermore, 31.1 per cent of the
cropped lesions touch the crop border, meaning that growth across that
border is censored without a missingness flag. A detailed treatment of
the spatial grid is deferred to~\S\ref{sec:data:spatial}.

For each eye in the cohort, the dataset comprises a sequence of
follow-up visits at irregular calendar intervals, each visit
contributing one eleven-channel en-face state. Patient-level
covariates --- in particular age and biological sex --- accompany the
imaging sequence; the operational details of how they are extracted
and encoded are described in~\S\ref{sec:data:covariates}. Because GA
is bilateral in approximately 39--50\,\% of patients at the time of
initial diagnosis with strong inter-eye correlation
\citep{Yehoshua2011,Boyer2017}, the dataset is split at the patient
level rather than the eye level, as motivated
in~\S\ref{sec:data:splits}.


\section{Partial Differential Equations and Numerical Solvers}
\label{sec:background:pde-solvers}

This section reviews the subset of classical PDE-solver theory that is
load-bearing for the remainder of the thesis. The exposition is
deliberately narrow: it covers temporal PDEs in conservation form, the
method-of-lines reduction to a system of ordinary differential
equations, finite-difference and finite-volume spatial stencils, the
elementary time integrators that close the loop on a discrete update
rule, and the distinction between uniform and adaptive meshes. The aim
is to provide the language in which the neural solvers
of~\S\ref{sec:background:neural-pde}--\S\ref{sec:background:mm-pde}
are defined; nothing further is required. In line with the project
convention, the two-dimensional shallow-water equation (SWE) serves as
the running concrete example throughout, both because the SWE features
prominently as a benchmark in the moving-mesh literature
\citep{Hu2024} and because its sharp wave fronts on a quadratic spatial
domain are visually similar to the active GA boundary on the
$49 \times 1024$ en-face grid.

\subsection{Temporal PDEs and conservation form}
\label{sec:background:pde-solvers:problem}

Following \citet{Brandstetter2022}, a temporal partial differential
equation in one time dimension and $D$ spatial dimensions is written as
\begin{equation}
  \partial_t u \;=\; F\!\left(t, x, u, \partial_x u, \partial_{xx} u, \dots\right),
  \qquad (t, x) \in [0, T] \times \mathcal{X},
  \label{eq:bg:temporal-pde}
\end{equation}
where $\mathcal{X} \subset \mathbb{R}^D$ is the spatial domain and the
solution $u: [0, T] \times \mathcal{X} \to \mathbb{R}^n$ takes values in
an $n$-dimensional state space. The problem is closed by an initial
condition $u(0, x) = u^{0}(x)$ and boundary conditions
$B[u](t, x) = 0$ on $[0, T] \times \partial \mathcal{X}$, where the
boundary operator $B$ may be of Dirichlet, Neumann, or periodic type.

A subclass of \eqref{eq:bg:temporal-pde} of particular interest in
this thesis is the family of equations expressible in
\emph{conservation form},
\begin{equation}
  \partial_t u \;+\; \nabla \cdot J(u) \;=\; 0,
  \label{eq:bg:conservation-form}
\end{equation}
where $J(u)$ is the flux of the conserved quantity $u$. The
conservation-form framing is the structural property that makes finite
volume methods \emph{exact} on the cell-integrated state and that
underpins the message-passing-as-stencil interpretation of neural
solvers in~\S\ref{sec:background:mp-pde}.

The SWE describe the depth-averaged motion of a thin layer of
incompressible fluid under gravity over a flat bottom and provide the
running example for this section.
% TODO: cite a canonical SWE reference (e.g. LeVeque 2002, "Finite
% Volume Methods for Hyperbolic Problems", Cambridge University Press,
% Chapter 13; or Vreugdenhil 1994, "Numerical Methods for Shallow-Water
% Flow", Springer) once added to references.bib -- the SWE are the
% running example for Section 2.2 and require a textbook citation.
With water column height $h(t, x, y)$ and depth-averaged horizontal
velocities $u(t, x, y)$ and $v(t, x, y)$, mass and momentum
conservation read
\begin{align}
  \partial_t h \;+\; \partial_x (h u) \;+\; \partial_y (h v) &= 0, \label{eq:bg:swe-mass}\\
  \partial_t (h u) \;+\; \partial_x \!\left( h u^{2} + \tfrac{1}{2} g h^{2} \right) \;+\; \partial_y \!\left( h u v \right) &= 0, \label{eq:bg:swe-momx}\\
  \partial_t (h v) \;+\; \partial_x \!\left( h u v \right) \;+\; \partial_y \!\left( h v^{2} + \tfrac{1}{2} g h^{2} \right) &= 0,
  \label{eq:bg:swe-momy}
\end{align}
where $g$ is gravitational acceleration. Equations
\eqref{eq:bg:swe-mass}--\eqref{eq:bg:swe-momy} are a coupled system of
hyperbolic conservation laws on the three-component state
$\mathbf{u} = (h, h u, h v)^{\top}$, and they fit
\eqref{eq:bg:conservation-form} with vector-valued flux
\[
  J(\mathbf{u}) \;=\;
  \begin{pmatrix}
    h u & h v \\
    h u^{2} + \tfrac{1}{2} g h^{2} & h u v \\
    h u v & h v^{2} + \tfrac{1}{2} g h^{2}
  \end{pmatrix}.
\]
The relevant features of this system for the present thesis are that the
solution is multi-component (foreshadowing the multi-channel state
tensor of~\S\ref{sec:method:multichannel}), that wave fronts can
develop sharp gradients in regions of small support
(foreshadowing the active GA boundary on a quasi-static interior and
exterior), and that a uniform discretisation of the spatial domain
wastes resolution in regions where the solution is smooth.

\subsection{The method of lines}
\label{sec:background:pde-solvers:mol}

The classical strategy for solving \eqref{eq:bg:temporal-pde}
numerically is the \emph{method of lines}, in which the spatial
variables are discretised first while the time variable is left
continuous \citep{Brandstetter2022}. The continuous spatial domain
$\mathcal{X}$ is replaced by a finite grid
$X = \{c_i\}_{i=1}^{N}$ of $N$ collocation points or cell centres, and
the value of the state at grid point $i$ and continuous time $t$ is
written $u_i(t)$. Spatial derivatives in $F(\cdot)$ are then
approximated by stencil-based operators evaluated on $X$, yielding the
semi-discrete system
\begin{equation}
  \frac{\mathrm{d} u_i}{\mathrm{d} t}
  \;=\; f_i\!\left(t, u_1(t), \dots, u_N(t)\right),
  \qquad i = 1, \dots, N,
  \label{eq:bg:mol-ode}
\end{equation}
in which only the time variable remains continuous. Equation
\eqref{eq:bg:mol-ode} is an ordinary differential equation in the
state vector $\mathbf{u}(t) = (u_1(t), \dots, u_N(t))^{\top}$ and can
be integrated with any standard ODE integrator. The full numerical
PDE solver is then a composition of two ingredients: a \emph{spatial
discretisation} that assembles the right-hand side of
\eqref{eq:bg:mol-ode}, and a \emph{time integrator} that advances
$\mathbf{u}(t)$ over a finite step $\Delta t$. The decomposition
matters because, as shown in~\S\ref{sec:background:mp-pde}, the
encode--process--decode neural solver of \citet{Brandstetter2022} can
be read directly as a learned non-linear instantiation of the same
two ingredients, with the GNN replacing the stencil and a residual
update replacing the ODE integrator.

\subsection{Spatial discretisation: stencils on a uniform grid}
\label{sec:background:pde-solvers:stencils}

For a uniform one-dimensional grid with spacing $\Delta x$, the
elementary stencil for the first spatial derivative is the
finite-difference approximation
\begin{equation}
  \left.\frac{\partial u}{\partial x}\right|_{x = x_i}
  \;\approx\; \frac{u_{i+1} - u_{i-1}}{2 \Delta x},
  \label{eq:bg:fdm-central}
\end{equation}
and analogous stencils exist for higher derivatives and higher
orders. More systematically, a polynomial of suitable degree is fitted
locally to a window of grid values, and the stencil coefficients are
read off as the corresponding derivative of the fitted polynomial
evaluated at the target grid point \citep{Brandstetter2022}. A
finite-difference method (FDM) is therefore a fixed linear operator
acting on the local neighbourhood of each grid point.

The finite-volume method (FVM) takes a different route, motivated by
the conservation form \eqref{eq:bg:conservation-form}. The spatial
domain is partitioned into control volumes (cells) $K_i$, and the
state value at cell $i$ is interpreted as the cell-averaged quantity
$\bar{u}_i = |K_i|^{-1} \int_{K_i} u \, \mathrm{d}x$. Integrating
\eqref{eq:bg:conservation-form} over $K_i$ and applying the
divergence theorem yields, in one spatial dimension,
\begin{equation}
  \frac{\mathrm{d} \bar{u}_i}{\mathrm{d} t}
  \;=\; \frac{1}{\Delta x_i}
  \!\left( J^{}_{i - 1/2} - J^{}_{i + 1/2} \right),
  \label{eq:bg:fvm-update}
\end{equation}
where $J^{}_{i \pm 1/2}$ are numerical flux estimates at the left and
right cell boundaries \citep{Brandstetter2022}. The cell average
$\bar{u}_i$ is exactly conserved up to boundary fluxes, which is the
discrete analogue of the local-conservation property of the underlying
PDE. Higher-order conservative reconstructions such as WENO5 add a
non-linear convex combination of sub-stencils that adapts to local
solution smoothness, allowing the scheme to capture sharp gradients
without introducing spurious oscillations \citep{Brandstetter2022}.

For the SWE in two dimensions, both FDM and FVM stencils generalise
in the natural way: the partial derivatives in
\eqref{eq:bg:swe-mass}--\eqref{eq:bg:swe-momy} are replaced by
finite-difference stencils on a tensor-product grid in the FDM case,
or by surface-flux estimates on the boundaries of two-dimensional
control volumes in the FVM case. In both cases the spatial
discretisation is local: the right-hand side of
\eqref{eq:bg:mol-ode} at grid point $i$ depends only on a small
neighbourhood of $i$. Locality is the structural property that
message-passing graph neural networks exploit explicitly when they
replace classical stencils, and it is the property that the dynamics
of GA --- atrophy concentrated at the lesion boundary, with the
lesion interior and the surrounding healthy retina both quasi-static
\citep{Yehoshua2011} --- already satisfies on the en-face grid.

\subsection{Time integration}
\label{sec:background:pde-solvers:time}

Once the spatial discretisation has reduced \eqref{eq:bg:temporal-pde}
to the system of ordinary differential equations
\eqref{eq:bg:mol-ode}, any standard ODE integrator closes the loop. The
simplest choice is the explicit Euler step
\begin{equation}
  u_i^{k + 1} \;=\; u_i^{k}
  \;+\; \Delta t \, f_i\!\left(t_k, u_1^{k}, \dots, u_N^{k}\right),
  \label{eq:bg:euler}
\end{equation}
with $u_i^{k} \approx u_i(t_k)$ and $\Delta t = t_{k+1} - t_k$. Higher
order is reached by replacing the single-stage Euler update with a
Runge--Kutta scheme; the fourth-order Runge--Kutta integrator (RK4) is
the standard choice for benchmark generation in the neural-PDE-solver
literature, used for instance in combination with WENO5 spatial
reconstruction by \citet{Brandstetter2022} to produce ground-truth
trajectories.

For the running SWE example, an explicit time-stepping scheme advances
the discretised state $\mathbf{u}^{k} = (h^{k}, (hu)^{k}, (hv)^{k})$ by
evaluating the discretised flux divergence on $\mathbf{u}^{k}$ and
applying \eqref{eq:bg:euler} or its Runge--Kutta extension. Stability
of the scheme is governed by a Courant--Friedrichs--Lewy (CFL)
condition that constrains $\Delta t$ in proportion to the spatial
spacing $\Delta x$ and the maximum wave speed
$\sqrt{g h} + \max(|u|, |v|)$; on a finer mesh, a smaller $\Delta t$
is required.

The form of \eqref{eq:bg:euler} is the structural template for the
\emph{residual} decoder of MP-PDE \citep{Brandstetter2022}, in which
the right-hand side $f_i$ is replaced by the output of a graph neural
network. The thesis adapts that template to the irregular
$\Delta t$ of clinical visits, replacing the fixed step size by the
per-window observation interval; the formal definition is given
in~\S\ref{sec:method:dt-residual}.

\subsection{Uniform and adaptive meshes}
\label{sec:background:pde-solvers:mesh}

The spatial discretisations described above all rest on a fixed
\emph{mesh} that partitions $\mathcal{X}$ into grid points or cells.
The simplest choice is a \emph{uniform mesh}, in which all cells have
the same size and the same shape; for a rectangular domain in two
spatial dimensions, this is a tensor-product grid with constant
spacings $\Delta x$ and $\Delta y$. A uniform mesh is the natural
discretisation when the solution is approximately equally hard to
resolve everywhere on the domain, which is rarely the case for
problems with localised structure.

When the solution exhibits sharp gradients in some regions and is
smooth in others --- as the SWE does in the vicinity of a wave front,
and as the GA state does in the vicinity of the active lesion boundary
--- a uniform mesh either over-resolves the smooth regions or
under-resolves the sharp ones. Two families of mesh-adaptation
strategies have been developed in classical numerical analysis to
address this trade-off \citep{Hu2024}. \emph{$h$-adaptation}
adds and removes vertices to refine or coarsen the mesh locally; the
total number of degrees of freedom changes over time and the mesh
topology is dynamic. \emph{$r$-adaptation}, sometimes called the
\emph{moving-mesh} approach, keeps the number of vertices and the
graph topology fixed and only \emph{moves} existing vertices toward
regions where the solution is hardest to approximate. The latter
strategy preserves the dimensionality and connectivity of the discrete
state across time steps, which is the property that makes it
compatible with graph-neural-network solvers; \S\ref{sec:background:mm-pde}
develops this point in detail.

The classical $r$-adaptation problem is well-posed only after a
\emph{monitor function} $M(x)$ has been specified, encoding where on
$\mathcal{X}$ the mesh should refine; the moving-mesh map then
redistributes vertex positions so that each cell carries the same
amount of monitor-weighted volume \citep{Hu2024}. \S\ref{sec:background:mm-pde}
shows how \citet{Hu2024} reformulate this redistribution as a
single-scalar Monge--Amp\`ere equation that can be solved without
ground-truth meshes, and how the resulting mesh mover plugs into a
neural PDE solver. The point for the present section is structural:
classical numerical analysis already recognises that fixed uniform
meshes are wasteful for problems with localised dynamics, and the
moving-mesh inductive bias of MM-PDE is a learned analogue of an
established classical idea --- one that on the GA imaging substrate is
motivated by the empirical fact that lesion progression is
concentrated at a thin boundary and is direction-dependent
\citep{Yehoshua2011,Singh2025}, while the lesion interior and the
surrounding healthy retina are quasi-static.

The remainder of Chapter~\ref{ch:background} traces the same chain of
ingredients --- discrete state on a mesh, local stencil, residual
update, optionally adapted mesh --- through their neural counterparts.


\section{Neural PDE Solvers and Surrogate Architectures}
\label{sec:background:neural-pde}

Neural PDE solvers broadly separate into two paradigms. Neural
operators learn a continuous map
$M: [0, T] \times \mathcal{F} \to \mathcal{F}$ such that
$M(t, u_0) = u(t)$, thereby predicting the solution at any arbitrary
time directly from the initial condition. By contrast, autoregressive
solvers learn a local update operator
$u(t + \Delta t) = A(\Delta t, u(t))$ and are applied iteratively to
unroll a trajectory. Note that when the step size $\Delta t$ is fixed
by the dataset, the explicit dependence on it is usually suppressed in
the notation. This thesis cannot suppress it because clinical
observation intervals are fundamentally irregular.

Neural operators take several structural forms. Finite-dimensional
variants are bound to the specific grid geometry they were trained on.
Infinite-dimensional variants decouple the learned operator from the
spatial discretisation altogether.
% TODO: cite Li2021 (Li, Kovachki, Azizzadenesheli, Liu, Bhattacharya,
% Stuart & Anandkumar, "Fourier Neural Operator for Parametric Partial
% Differential Equations", ICLR 2021)
The Fourier Neural Operator (Li et al., 2021) parameterises a kernel
integral operator in the Fourier domain. This formulation gives the
network global spatial support in a single layer.
% TODO: cite Lu2021 (Lu, Jin, Pang, Zhang & Karniadakis, "Learning
% nonlinear operators via DeepONet based on the universal approximation
% theorem of operators", Nature Machine Intelligence 3(3):218-229, 2021)
Alternatively, DeepONet (Lu et al., 2021) factorises the operator into
a branch network that encodes the input function and a trunk network
that encodes the continuous query coordinate. A standard limitation
applies to both variants: the trained operator is inherently tied to
the equation family it was trained on, and extrapolation beyond the
temporal horizon seen during training is reliably poor.

The autoregressive framing is the natural fit for clinical
longitudinal data. Each visit-to-visit transition fundamentally
constitutes one step. Furthermore, there is no shared temporal origin
($t = 0$) across patients, making a direct time-dependent mapping
ill-posed. Finally, the required prediction horizon is set by the
clinic rather than by the model, which necessitates a system that can
be stepped forward as many times as a given patient's schedule
dictates.

The principal cost of the autoregressive framing is error
accumulation. After the first step, the solver consumes its own output
rather than numerical ground truth. The input distribution at test
time therefore differs systematically from the distribution seen
during training. This distribution shift is the central training
difficulty for autoregressive rollouts.
\S\ref{sec:background:mp-pde} describes the two primary mechanisms
proposed to mitigate it in the graph-network literature.

Beyond graph networks, several distinct architecture families are used
as PDE surrogates.
% TODO: cite Ronneberger2015 (Ronneberger, Fischer & Brox, "U-Net:
% Convolutional Networks for Biomedical Image Segmentation",
% MICCAI 2015, pp. 234-241)
% TODO: cite Gupta2023 (Gupta & Brandstetter, "Towards
% Multi-spatiotemporal-scale Generalized PDE Modeling", TMLR 2023)
The U-Net (Ronneberger et al., 2015) is an encoder-decoder
architecture with skip connections and a multi-scale receptive field.
Although originally introduced for biomedical image segmentation, it
serves as a standard strong surrogate baseline in the neural-PDE
literature (Gupta \& Brandstetter, 2023).
% TODO: cite LiuSchiaffini2024 (Liu-Schiaffini, Berner, Bonev, Kurth,
% Azizzadenesheli & Anandkumar, "Neural Operators with Localized
% Integral and Differential Kernels", ICML 2024)
Surrogate architectures also vary fundamentally in their spatial
context. Global spectral operators, such as the Fourier Neural
Operator, contrast directly with local convolutional or
message-passing operators. Hybrid architectures combine these
properties by adding a small local kernel bypass to a global operator
to capture high-frequency details more effectively
(Liu-Schiaffini et al., 2024).
% TODO: cite Gao2019 (Gao & Ji, "Graph U-Nets", ICML 2019)
Graph U-Nets (Gao \& Ji, 2019) provide the corresponding multi-scale
counterpart for data structured as graphs.
% TODO: cite Lienen2022 (Lienen & Gunnemann, "Learning the Dynamics of
% Physical Systems from Sparse Observations with Finite Element
% Networks", ICLR 2022)
Finite Element Networks (Lienen \& G\"unnemann, 2022) adopt a different
numerical perspective, building the neural update directly from a
learned finite-element discretisation and carrying an explicit
transport (advection) term.
% TODO: cite Chen2018 (Chen, Rubanova, Bettencourt & Duvenaud, "Neural
% Ordinary Differential Equations", NeurIPS 2018)
% TODO: cite Ott2021 (Ott, Katiyar, Hennig & Tiemann, "ResNet After
% All: Neural ODEs and Their Numerical Solution", ICLR 2021)
% TODO: cite Krishnapriyan2023 (Krishnapriyan, Queiruga, Erichson &
% Mahoney, "Learning continuous models for continuous physics",
% Communications Physics 6:319, 2023)
Finally, because a residual step resembles one explicit Euler step, a
Neural-ODE reading is sometimes applied (Chen et al., 2018). Under
this interpretation, a higher-order integrator such as a Runge--Kutta
scheme can in principle replace the residual step. It must be stated
plainly that this reading is only legitimate if the trained model is
nearly invariant to the solver used at inference. This property must
be tested empirically rather than assumed, as classical solver
guarantees do not automatically transfer to neural parameterisations
(Ott et al., 2021; Krishnapriyan et al., 2023).

These architecture families differ in a small number of structural
properties. The primary distinctions lie in whether the operator has
spatial context at all, how far that context physically reaches,
whether it combines global and local information, and whether it
carries an explicit transport term. This thesis refers to these
structural differences as inductive-bias ingredients.
Chapter~\ref{ch:experiments} isolates and measures them.


\section{MP-PDE: Message-Passing Neural PDE Solvers}
\label{sec:background:mp-pde}

\citet{Brandstetter2022} introduced an autoregressive neural solver for
time-dependent PDEs in conservation form. The solver operates on a
spatial graph, leveraging a direct analogy where grid points or
finite-volume cells become nodes, and numerical stencils become edges.
The spatial graph is constructed dynamically via a cutoff radius on
regular grids, or by connecting $k$-nearest-neighbours on irregular
meshes.

% TODO: cite Battaglia2018 (Battaglia et al., "Relational inductive
% biases, deep learning, and graph networks", arXiv 2018)
% TODO: cite SanchezGonzalez2020 (Sanchez-Gonzalez et al., "Learning to
% simulate complex physics with graph networks", ICML 2020)
The architecture follows a standard encode--process--decode pattern
(Battaglia et al., 2018; Sanchez-Gonzalez et al., 2020). An encoder,
implemented as a per-node two-layer multilayer perceptron (MLP), lifts
the initial node features into a latent space. The encoder input takes
the form
$f_i^0 = \epsilon([u_i^{k-K:k}, x_i, t_k, \theta_{\text{PDE}}])$.

The processor applies multiple rounds of message passing to update the
latent node state; the original paper sets the number of rounds to six.
The message function is defined as
$m_{ij}^m = \phi(f_i^m, f_j^m, u_i^{k-K:k} - u_j^{k-K:k}, x_i - x_j,
\theta_{\text{PDE}})$, and the subsequent node update is
$f_i^{m+1} = \psi(f_i^m, \sum_j m_{ij}^m, \theta_{\text{PDE}})$. Both
$\phi$ and $\psi$ are parameterised as two-layer MLPs with skip
connections and normalisation layers. The inclusion of relative
positions $x_i - x_j$ directly encodes translational symmetry into the
network. The solution differences $u_i - u_j$ act as a learned local
difference operator, establishing a flexible learned stencil.

A core representational claim of the MP-PDE framework is that message
passing can approximate classical numerical schemes. The original
authors mapped roughly one message-passing layer to a finite-difference
scheme, two layers to finite volumes, and three layers to a WENO5
reconstruction.

The equation feature vector, denoted $\theta_{\text{PDE}}$, holds the
PDE coefficients and the boundary-condition identifiers. Crucially,
this vector is injected both at the encoder and at every single
message-passing layer. This persistent injection is what enables a
single trained solver to generalise across a family of equations.
Ablating the $\theta_{\text{PDE}}$ injection degrades performance most
severely when the equation parameters vary across the dataset.

The decoder is a shallow one-dimensional convolutional neural network
with weights shared across the spatial dimensions. It treats the final
hidden vector of each node as a temporal signal and convolves along it,
producing a sequence of updates $d_i = (d_i^1, \dots, d_i^K)$. A residual update
$u_i^{k+l} = u_i^k + (t_{k+l} - t_k) d_i^l$ produces the final
prediction for each substep $1 \leq l \leq K$. This formulation is
motivated by numerical consistency: the update must necessarily
converge to the true solution as the temporal step size approaches
zero.

Two distinct mechanisms mitigate the accumulation of rollout errors.
% TODO: verify the bundle sizes against Brandstetter et al. (2022); the
% set {20, 25, 50} comes from the code-side notes and could not be located
% in the paper's text. Remove the parenthesis if it cannot be confirmed.
The first is temporal bundling. The model predicts $K$ future steps per
forward pass (with $K \in \{20, 25, 50\}$ in the original evaluation).
This reduces the total number of solver calls required along a given
rollout by a factor of $K$, which in turn reduces the number of
accumulated distribution shifts the model is exposed to.

The second mechanism is the pushforward trick, an adversarial stability
loss. During training, the input state is perturbed by the model's own
one-step prediction from the preceding state. This ensures that the
perturbation applied to the training input is drawn from the exact
distribution the solver will subsequently face at inference time. The
trick is implemented by unrolling two steps and backpropagating only
through the last; gradients on the earlier step are deliberately cut.
Cutting the gradients is both faster and more stable, because
backpropagating through the perturbation itself would teach the model
to make the perturbation small rather than learning to correct it. This
technique empirically outperforms training with Gaussian noise, which
improves stability but consistently incurs a cost to predictive
accuracy.

The pushforward trick is theoretically grounded in the concept of zero
stability. An update operator is considered zero-stable if
$\lVert A(u^0 + \epsilon) - u^1 \rVert \leq \kappa \lVert \epsilon
\rVert$ for a small constant $\kappa$. The pushforward loss directly
minimises this constant $\kappa$. Zero stability, together with
consistency, provides the fundamental convergence argument that the
MP-PDE architecture relies upon.

At scale, the original MP-PDE models comprise roughly one million
parameters. The hidden representation size is of the order of 128, and
the networks are trained using the AdamW optimiser.

The original paper states three clear limitations. High-quality
numerical ground truth is strictly required for training. There are no
formal accuracy guarantees. Finally, a uniform mesh is wasteful in
smooth regions and under-resolved at sharp gradients. This last
limitation directly motivates the moving-mesh extension detailed in the
next section.

This thesis keeps the residual update formulation and the pushforward
trick. It modifies the architecture by conditioning the update on a
variable $\Delta t$ and removes temporal bundling entirely. The
specifics of this adaptation are deferred to
Chapter~\ref{ch:method}.


\section{MM-PDE: Moving Mesh PDE Solvers}
\label{sec:background:mm-pde}

\citet{Hu2024} extend message-passing neural solvers by integrating an
adaptive moving mesh. A uniform mesh is fundamentally inefficient for
problems that exhibit localised dynamics. It wastes computational
resolution in smooth regions while simultaneously under-resolving sharp
gradients, such as wave fronts.

Mesh adaptation strategies separate into two families. The
$h$-adaptation approach adds and removes vertices to dynamically refine
or coarsen the mesh. This changes both the topology and the total
degrees of freedom between steps, a property that is hostile to
standard graph-network training. Conversely, $r$-adaptation, or
moving-mesh adaptation, moves a fixed set of nodes while maintaining
both the node count and the original graph topology. Prior neural
$r$-adaptation methods were entirely supervised and required
ground-truth optimal meshes generated by an expensive classical solver.

The moving-mesh framework is governed by a monitor function. The
monitor is defined as $M = (1 + \lVert \nabla u \rVert_2 / \alpha) I$,
a matrix-valued function that takes large values precisely where the
state solution varies sharply. The identity factor $I$ ensures that the
refinement is isotropic. The addition of a constant $+1$ keeps $M$
bounded away from zero, which prevents nodes from collapsing entirely
in smooth regions. The global scaling constant $\alpha$ dictates the
trade-off between mesh adaptivity and numerical stability. The
corresponding mesh density is given by $\rho(x) = \sqrt{\det M(x)}$.
This specific functional form is derived from an established upper
bound on finite-element interpolation error.

The equidistribution principle formalises the objective of the mesh
adaptation. A coordinate map $f$ produces an $M$-uniform mesh if the
integral of the density $\rho$ over the image of every single cell is
equal. Specifically, this integral must equal $\sigma / N$, where
$\sigma$ is the total integral of $\rho$ over the domain and $N$ is the
total number of cells. Intuitively, every cell must carry the identical
amount of monitor-weighted information, which forces cells to shrink in
regions where $M$ is large.

Because infinitely many maps satisfy the equidistribution principle,
choosing the optimal map is framed as an optimal transport problem. The
objective is to select the mapping closest to the identity. By
Brenier's theorem, the unique solution to this problem is the gradient
of a convex scalar potential, such that $f = \nabla \phi$. Substituting
this gradient into the equidistribution constraint yields the
Monge--Amp\`ere equation,
\[
  \det H(\phi) \;=\; \frac{\sigma}{|\Omega| \, \rho(\nabla \phi)},
\]
which is a fully non-linear second-order PDE for the scalar potential
$\phi$. The problem is constrained by the boundary condition that the
domain boundary maps exactly to itself. A scalar potential is learned
rather than the full vector map because the convexity of $\phi$ is
mathematically equivalent to a non-singular Jacobian of $f$. A
non-singular Jacobian guarantees that the resulting mesh does not
tangle. Crucially, enforcing convexity on a scalar potential is far
easier to regularise than imposing a global non-tangling constraint on
a vector field.

The network architecture reformulates the scalar potential residually
as $\phi(x) = 0.5 \|x\|^2 + \psi(x)$. The gradient therefore becomes
$\nabla \phi = x + \nabla \psi$. The network is tasked with predicting
only the displacement field $\psi$. This formulation removes the
trivial identity component from the network's objective and ensures
that the coordinate map is approximately the identity mapping at
initialisation.

The Data-free Mesh Mover (DMM) follows a DeepONet-style architecture. A
branch network encodes the full spatial state $u$, and a trunk network
processes a continuous query coordinate. Together, they output the
scalar potential. The DMM is trained entirely without mesh supervision
via a dedicated physics loss
$L = L_{\text{eq}} + \beta L_{\text{bound}} + \gamma L_{\text{convex}}$.
In this formulation, $L_{\text{eq}}$ is the squared residual of the
Monge--Amp\`ere equation evaluated at sampled interior points,
$L_{\text{bound}}$ enforces the boundary condition, and
$L_{\text{convex}}$ penalises violations of convexity. The convexity
term is critical; a non-convex potential inevitably yields a tangled
and useless mesh, regardless of how small the equation residual
becomes. Training uses the Adam optimiser followed by L-BFGS
fine-tuning applied exclusively to the last layer. The DMM is
pretrained and subsequently frozen during solver training. This
decouples a well-posed geometric problem from the main solver objective
and allows the identical mesh mover to be reused across different
experiments.

The full moving-mesh update operator is constructed as a dual-branch
composition, defined as $M(u_t) = G_1(u_t) + I(G_2, \tilde{f}, u_t)$.
Here, $G_1$ represents a graph neural network operating on the uniform
mesh, $G_2$ is a graph neural network operating on the moved mesh,
$\tilde{f}$ is the coordinate map predicted by the DMM, and $I$ is a
learned interpolation operator mapping features back to uniform
coordinates. The uniform branch supplies a stable global view of the
domain, while the moved branch supplies a locally refined correction
where gradients are steep.

The interpolation network, ItpNet, learns adaptive interpolation
weights directly from the moved-mesh geometry rather than relying on
fixed finite-element basis functions. It is pretrained with its own
optimiser before joint dual-branch training begins, because an
untrained interpolation would feed unstructured noise into the uniform
branch, destabilising the system.

% TODO: cite HuangRussell2011 (Huang & Russell, "Adaptive Moving Mesh
% Methods", Springer, Applied Mathematical Sciences vol. 174, 2011)
The classical background for adaptive moving meshes on which this
framework builds is covered extensively by Huang \& Russell (2011).
\citet{Hu2024} reported several architectural ablations that confirm
the dual-branch necessity. Removing the moved branch degrades
performance entirely to the uniform-mesh baseline. Removing the uniform
branch degrades performance even further, because the moved mesh
inherently under-covers smooth regions. Replacing the DMM-generated
mesh with a random mesh isolates and degrades the performance of the
moved branch specifically. Finally, removing the interpolation
pretraining severely harms overall training stability, even if it does
not permanently cap the final accuracy.

Whether mesh adaptation provides a measurable benefit for slowly
evolving clinical data remains an open question, and one that this
thesis tests. How a monitor function should be defined over a
multi-channel state tensor is similarly treated as an open question
deferred to Chapter~\ref{ch:method}.


\section{Related Work}
\label{sec:background:related-work}

Deep learning models for OCT-based age-related macular degeneration and
geographic atrophy progression divide broadly into cohort-level risk
models and spatial-progression models.
% TODO: cite SchmidtErfurth2018 (Schmidt-Erfurth et al., "Prediction of
% individual disease conversion in early AMD using artificial
% intelligence", IOVS 59:3199-3208, 2018)
Schmidt-Erfurth et al. (2018) reported on the prediction of individual
disease conversion in early AMD using artificial intelligence.
% TODO: cite Bogunovic2017 (Bogunovic et al., "Machine learning of the
% progression of intermediate AMD based on OCT imaging",
% IOVS 58:BIO141-BIO150, 2017)
Bogunovic et al. (2017) demonstrated the machine learning of
intermediate-AMD progression directly from OCT imaging. For established
geographic atrophy, \citet{Vogl2021} provided a topographic progression
analysis. The broader landscape of deep-learning progression models in
this domain is reviewed extensively by \citet{Lad2023} and
\citet{Singh2025}. \citet{Vallino2024} explicitly argued that
artificial intelligence applied to OCT can provide predictive
probabilities of GA development over time, cementing OCT as a
predictive substrate for structural endpoints.

% TODO: cite Mai2024 (Mai, Lachinov, Reiter, Riedl, Grechenig,
% Bogunovic & Schmidt-Erfurth, "Deep Learning-Based Prediction of
% Individual Geographic Atrophy Progression from a Single Baseline
% OCT", Ophthalmology Science 4(4):100466, 2024)
The closest prior work on the specific prediction task addressed here
is by Mai et al. (2024), who predicted individual GA progression from a
single baseline OCT on the exact MUW cohort used in this thesis. The
input regime, however, is fundamentally different. While that prior
work operates directly on raw OCT volumes, the present thesis consumes
pre-segmented lesion masks and layer surfaces. This difference must be
kept in mind whenever the approaches are compared.
% TODO: cite Salvi2025 (Salvi et al., "Deep Learning to Predict the
% Future Growth of Geographic Atrophy from Fundus Autofluorescence",
% Ophthalmology Science 5(2):100635, 2025)
As a precedent for dense convolutional architectures on this task,
Salvi et al. (2025) predicted future GA growth using a two-dimensional
U-Net on fundus autofluorescence images, establishing a relevant
multi-scale spatial benchmark on an alternative modality.

The prevailing formulation in the clinical literature rests on
cohort-level, categorical, time-to-conversion prediction. By contrast,
the formulation used here dictates a per-eye, per-time-step prediction
of the full spatial next state. Consequently, this thesis sits at the
intersection of applied medical-imaging deep learning and PDE-solver
inductive biases.
% TODO: check the literature for any prior application of neural PDE
% solvers to biomedical disease progression or organ modelling. Add a
% citation if one exists; otherwise note the absence explicitly.
The application of neural PDE solvers to biomedical problems remains
sparse.


%%
%%%%%%%%%%%%%%%%%%%%%%%%%%%%%%%%%%%%%%%%%%%%%%%%%%%%%%%%%%%%%%%%%%%%%%%%%%%%%%%%
